\chapter{Heart Valve Replacement}

\section*{What Is This Surgery?}
Heart valve replacement is a definitive cardiac surgical procedure involving the meticulous removal of a severely diseased or dysfunctional native heart valve and its precise substitution with a prosthetic valve. This complex intervention is primarily undertaken to address critical conditions such as severe valvular stenosis (pathological narrowing) or regurgitation (incompetent leakage) predominantly affecting the aortic or mitral valves, although tricuspid and pulmonary valves can also be replaced in specific, less common clinical scenarios. The overarching goal of this surgery is to restore normal, unidirectional blood flow dynamics through the cardiac chambers, thereby alleviating debilitating symptoms of heart failure, improving compromised cardiac function, and ultimately enhancing patient survival and long-term quality of life. The critical decision between implanting a mechanical or a bioprosthetic (tissue) valve is carefully tailored to various patient-specific factors, including age, lifestyle considerations, and the individual's tolerance for lifelong anticoagulation therapy.

\section*{Alternative Names}
\begin{itemize}
  \item Surgical Valve Replacement
  \item Aortic Valve Replacement (AVR)
  \item Mitral Valve Replacement (MVR)
  \item Tricuspid Valve Replacement (TVR)
  \item Pulmonary Valve Replacement (PVR)
  \item Open Heart Valve Surgery
  \item Valve Replacement Surgery
\end{itemize}

\section*{Relevant Surgical Anatomy}
Successful heart valve replacement necessitates an intimate understanding of the intricate anatomy of the heart and its surrounding structures. The heart is a four-chambered muscular pump, comprising two atria and two ventricles, separated by a septum. Blood flow is regulated by four valves: the tricuspid and mitral (atrioventricular) valves, and the pulmonary and aortic (semilunar) valves. Each valve consists of delicate leaflets or cusps anchored to a fibrous annulus, which forms the structural support. The mitral valve, located between the left atrium and left ventricle, typically has two leaflets (anterior and posterior) connected by chordae tendineae to papillary muscles within the left ventricle. The aortic valve, situated between the left ventricle and the aorta, usually has three semilunar cusps (right, left, and non-coronary) that open into the ascending aorta, immediately distal to which lie the ostia of the coronary arteries.

The proximity of the cardiac conduction system is paramount during valve surgery. The atrioventricular (AV) node and the Bundle of His lie in close relation to the fibrous trigones and the aortic and mitral annuli, making them vulnerable to injury during valve excision and suture placement, potentially leading to heart block. The great vessels, including the ascending aorta, pulmonary artery, and the superior and inferior vena cavae, are crucial for cannulation during cardiopulmonary bypass. The pericardium, a double-layered sac enclosing the heart, is incised to access the heart and then typically closed at the end of the procedure. Understanding the precise spatial relationships of these structures, along with their arterial supply (coronary arteries), venous drainage (cardiac veins), and lymphatic pathways, is fundamental to minimizing complications and achieving optimal surgical outcomes.

\section{Surgical Principle \& Rationale}
The fundamental principle of heart valve replacement is to correct severe valvular dysfunction that significantly impairs cardiac output and leads to debilitating symptoms of heart failure or irreversible myocardial damage. When a heart valve becomes stenotic, its narrowed orifice obstructs forward blood flow, forcing the heart to exert greater pressure to eject blood. This increased workload leads to compensatory myocardial hypertrophy, which, if uncorrected, eventually progresses to ventricular dilation and overt heart failure. Conversely, a regurgitant valve fails to close completely, allowing a portion of the blood to flow backward into the preceding chamber. This volume overload increases the workload on the cardiac chambers and reduces the effective forward flow to the systemic circulation. Both conditions compromise the heart's pumping efficiency and can be life-threatening if left untreated.

The surgical intervention aims to meticulously remove the diseased native valve leaflets and annulus, which are often pathologically calcified, fibrotic, or structurally compromised due to various etiologies. Following excision, a new, functional prosthetic valve is precisely implanted. This prosthetic valve, whether mechanical or bioprosthetic, is engineered to ensure competent, unidirectional blood flow, thereby normalizing intracardiac pressures, reducing the excessive cardiac workload, and restoring adequate systemic circulation. Key technical goals include achieving a secure, leak-free anastomosis between the prosthesis and the native annulus, optimizing valve size to ensure excellent hemodynamics without patient-prosthesis mismatch, and carefully preserving surrounding cardiac structures, such as the mitral subvalvular apparatus (chordae and papillary muscles), where possible, to maintain ventricular geometry and function.

\section{History \& Surgical Evolution}
The journey towards successful heart valve replacement represents one of the most significant triumphs in modern cardiac surgery, commencing in the mid-20th century. A pivotal prerequisite for this advancement was the development of cardiopulmonary bypass (the heart-lung machine) in the early 1950s by Dr. John Gibbon, which allowed surgeons to operate on a still, bloodless heart. Building upon this, in 1960, Dr. Albert Starr and Lowell Edwards achieved a monumental breakthrough by developing and successfully implanting the first clinically viable caged-ball valve prosthesis in a human patient. This pioneering work at the University of Oregon Medical School revolutionized the treatment of valvular heart disease and paved the way for the rapid development of other mechanical valve designs, including tilting disc valves (e.g., Bjork-Shiley) and bileaflet valves (e.g., St. Jude Medical), which offered improved hemodynamics and reduced thrombogenicity.

The 1970s witnessed the introduction of bioprosthetic (tissue) valves, derived from porcine aortic valves or bovine pericardium. These valves provided a crucial alternative for patients who could not tolerate lifelong anticoagulation, despite their limited durability compared to mechanical valves. Further evolution in the 21st century brought forth minimally invasive surgical approaches, such as ministernotomy or thoracotomy, aiming to reduce surgical trauma and accelerate recovery. Most notably, the landscape of valve replacement was dramatically transformed in 2002 with the first human transcatheter aortic valve implantation (TAVI, or TAVR in the US) performed by Dr. Alain Cribier. TAVR revolutionized the treatment of high-risk and inoperable patients with severe aortic stenosis, allowing valve replacement without open-heart surgery, and its indications continue to expand, representing a paradigm shift in cardiovascular care.

\section{Clinical Indications}
Heart valve replacement is indicated for patients with severe valvular heart disease who are symptomatic or show objective evidence of progressive cardiac dysfunction or adverse remodeling, even if asymptomatic. Specific clinical indications, often guided by international society guidelines (e.g., ACC/AHA, ESC), include:
\begin{itemize}
  \item Severe symptomatic aortic stenosis, characterized by classic symptoms such as angina pectoris, syncope, or dyspnea on exertion.
  \item Severe asymptomatic aortic stenosis with objective evidence of left ventricular dysfunction (ejection fraction <50\%), rapid progression, or exercise-induced symptoms.
  \item Severe symptomatic aortic regurgitation with left ventricular dilation or dysfunction, or severe asymptomatic aortic regurgitation with progressive left ventricular enlargement.
  \item Severe symptomatic mitral stenosis, typically presenting with dyspnea, pulmonary hypertension, or recurrent embolic events despite medical therapy.
  \item Severe symptomatic primary (degenerative) mitral regurgitation with left ventricular dysfunction or dilation, or severe asymptomatic primary mitral regurgitation with new-onset atrial fibrillation or pulmonary hypertension.
  \item Severe secondary (functional) mitral regurgitation that remains refractory to optimal medical therapy and cardiac resynchronization therapy in carefully selected patients, often considered in conjunction with coronary artery bypass grafting.
  \item Infective endocarditis causing severe valvular destruction, persistent infection despite appropriate antibiotic therapy, recurrent systemic emboli, or heart failure due to valvular incompetence.
  \item Valvular dysfunction resulting from congenital heart disease, traumatic injury, or other acquired conditions causing significant hemodynamic compromise and symptoms.
\end{itemize}

\section{Clinical Positioning}
For the vast majority of patients presenting with severe, symptomatic valvular heart disease, heart valve replacement is considered a primary and definitive treatment. It directly addresses the underlying structural pathology of the valve, aiming to restore normal cardiac function, alleviate symptoms, and prevent irreversible myocardial damage. In cases of severe aortic stenosis or regurgitation, or severe mitral stenosis, surgical valve replacement (or transcatheter aortic valve replacement for aortic stenosis in eligible patients) is often the first-line intervention once medical management is no longer sufficient to control symptoms or when specific guideline-based criteria for intervention are met. The decision to proceed is typically made by a multidisciplinary Heart Team, weighing the risks and benefits of surgery versus medical management or alternative less invasive procedures.

However, in certain complex clinical scenarios, valve replacement may be positioned as a secondary, salvage, or adjunct procedure. For instance, in severe secondary (functional) mitral regurgitation, valve replacement (or repair) may be considered after optimal medical therapy, including guideline-directed heart failure treatment and potentially cardiac resynchronization therapy, has failed to improve symptoms or left ventricular remodeling. Similarly, valve replacement may be performed as an adjunct to other cardiac surgeries, such as coronary artery bypass grafting (CABG), when significant valvular disease coexists with coronary artery disease. The role of valve replacement as primary versus secondary intervention is meticulously guided by patient symptoms, left ventricular function, the specific characteristics and severity of the valvular lesion, and the patient's overall surgical risk profile.

\section{Surgical Technique \& Procedure}
Heart valve replacement is a complex open-heart procedure requiring meticulous surgical technique. The procedure typically commences with the administration of general anesthesia, followed by careful patient positioning, usually supine, with arms tucked to facilitate sternal access. A median sternotomy incision is then made, providing direct and comprehensive access to the heart and great vessels.

The patient is subsequently placed on cardiopulmonary bypass (CPB), commonly referred to as the heart-lung machine. This critical step involves cannulation of the aorta and vena cavae, diverting systemic venous blood to the CPB circuit for oxygenation and carbon dioxide removal, and returning oxygenated blood to the aorta. This allows the surgeon to operate on a still, bloodless heart. Myocardial protection is achieved by administering a cold cardioplegic solution into the coronary arteries, which arrests the heart and minimizes ischemic damage during the period of aortic cross-clamping.

Key operative steps for surgical valve replacement include:
\begin{enumerate}
  \item  \textbf{Cardiopulmonary Bypass Initiation \& Aortic Cross-clamping:} After systemic heparinization, the patient is cannulated, and CPB is initiated. The ascending aorta is then cross-clamped to isolate the heart from systemic circulation, and cardioplegia is delivered to achieve cardiac arrest.
  \item  \textbf{Valve Exposure:} For aortic valve replacement, an aortotomy (incision in the aorta) is made just above the aortic root. For mitral valve replacement, an incision is typically made in the left atrium (left atriotomy) or through the interatrial groove.
  \item  \textbf{Diseased Valve Excision \& Annular Debridement:} The diseased native valve leaflets or cusps are carefully excised using fine scissors and scalpels. Any calcification, fibrotic tissue, or vegetations are meticulously debrided from the valve annulus to create a clean, smooth bed for the new prosthesis.
  \item  \textbf{Valve Sizing \& Prosthesis Selection:} The native valve annulus is precisely measured using specialized sizers to determine the optimal size of the prosthetic valve. The choice between mechanical and bioprosthetic valves is made preoperatively, considering patient factors.
  \item  \textbf{Prosthesis Implantation:} The selected prosthetic valve is meticulously sutured into the native valve annulus using multiple non-absorbable sutures, typically in an interrupted or continuous fashion. Care is taken to ensure secure, even placement and a leak-free seal.
  \item  \textbf{De-airing \& Bypass Weaning:} Before removing the aortic cross-clamp, the heart chambers are carefully de-aired to prevent air embolism. The heart is then allowed to reperfuse, and the patient is gradually weaned off cardiopulmonary bypass, often with inotropic support to optimize cardiac function.
  \item  \textbf{Closure:} The aortotomy or atriotomy is closed, and the sternum is reapproximated with stainless steel wires. The chest wall and skin are then closed in layers. Mediastinal and/or pleural chest tubes are typically placed to drain fluid and blood from the chest cavity. Temporary epicardial pacing wires may also be placed for immediate postoperative rhythm management.
\end{enumerate}

\section*{Preoperative Workup \& Preparation}
Thorough preoperative evaluation and preparation are paramount to optimizing patient outcomes, minimizing surgical risks, and ensuring a smooth recovery following heart valve replacement. This comprehensive process typically involves:
\begin{itemize}
  \item \textbf{Comprehensive Medical History and Physical Examination:} To assess overall health status, identify significant comorbidities (e.g., diabetes, hypertension, renal insufficiency, pulmonary disease), and evaluate functional capacity.
  \item \textbf{Laboratory Tests:} Including a complete blood count, coagulation profile (PT/INR, PTT), renal and liver function tests, electrolytes, blood glucose, and a blood type and crossmatch for potential transfusions.
  \item \textbf{Cardiac Imaging:} Transthoracic echocardiography (TTE) and often transesophageal echocardiography (TEE) are essential to confirm valvular pathology, quantify its severity, assess ventricular function, and measure valve gradients. Coronary angiography is mandatory for patients over 40 years old, those with risk factors for coronary artery disease, or with symptoms suggestive of ischemia, to rule out concurrent coronary lesions requiring revascularization.
  \item \textbf{Dental Clearance:} A crucial step to identify and treat any active dental infections or significant periodontal disease, as these can serve as a source for postoperative prosthetic valve endocarditis.
  \item \textbf{Pulmonary Function Tests (PFTs):} Recommended for patients with a significant history of respiratory disease or symptoms, to assess lung capacity and optimize respiratory status preoperatively.
  \item \textbf{Anesthesia and Cardiac Clearance:} Evaluation by an anesthesiologist and cardiologist to assess surgical risk (e.g., using STS risk scores), optimize medical management, and discuss the perioperative plan.
  \item \textbf{Medication Adjustments:} Anticoagulants (e.g., warfarin, direct oral anticoagulants) are typically discontinued several days prior to surgery and may be bridged with intravenous heparin if the risk of thromboembolism is high. Antiplatelet agents (e.g., aspirin, clopidogrel) are also usually held.
  \item \textbf{NPO Status:} Patients are instructed to be nil per os (NPO) for at least 6-8 hours before surgery to prevent aspiration during induction of anesthesia.
  \item \textbf{Preoperative Antibiotics:} Administered intravenously within one hour prior to surgical incision to prevent surgical site infection and prosthetic valve endocarditis.
  \item \textbf{Informed Consent:} A detailed discussion with the patient and family regarding the procedure, its potential benefits, inherent risks, available alternatives, and the expected recovery trajectory.
\end{itemize}

\section*{Contraindications \& Precautions}
**Absolute Contraindications:**
\begin{itemize}
  \item Irreversible neurological damage or severe cognitive impairment that would preclude a meaningful recovery and acceptable quality of life post-surgery.
  \item Severe, uncorrectable comorbidities with an exceedingly high predicted operative mortality (e.g., end-stage liver disease with coagulopathy, severe pulmonary hypertension with refractory right heart failure, active metastatic malignancy with a very short life expectancy).
  \item Active systemic infection or sepsis that is not adequately controlled by antibiotics, as this dramatically increases the risk of prosthetic valve endocarditis and other severe postoperative infections.
  \item Inability to tolerate cardiopulmonary bypass for open surgical valve replacement due to severe vascular disease or other factors, although transcatheter options may be considered in such cases.
  \item Severe porcelain aorta (extensive circumferential calcification of the ascending aorta), which makes aortic clamping hazardous and significantly elevates the risk of perioperative stroke.
\end{itemize}

**Relative Contraindications \& Precautions:**
\begin{itemize}
  \item Advanced age and frailty, which may significantly increase surgical risk and prolong recovery; in these patients, less invasive options like transcatheter aortic valve replacement (TAVR) or transcatheter edge-to-edge repair (TEER) for mitral regurgitation may be preferred.
  \item Significant renal insufficiency, necessitating careful fluid management, avoidance of nephrotoxic drugs, and consideration of perioperative dialysis.
  \item Severe chronic obstructive pulmonary disease (COPD) or other significant respiratory compromise, which increases the risk of postoperative pulmonary complications such as pneumonia and prolonged ventilation.
  \item Previous cardiac surgery, which increases the technical complexity and risk of reoperation due to dense adhesions, altered anatomy, and potential for injury to previously grafted vessels.
  \item Coagulopathy or bleeding diathesis, requiring meticulous management of anticoagulation, hemostasis, and blood product administration.
  \item Patient refusal or a lack of complete understanding of the procedure, its implications, and the commitment to long-term follow-up, necessitating further counseling and education.
\end{itemize}

\section*{Complications \& Risks}
Heart valve replacement is a major surgical procedure carrying inherent potential risks, which can be broadly categorized into general surgical complications and those specific to cardiac valve surgery.
\begin{itemize}
  \item \textbf{General Surgical Risks:}
  \begin{itemize}
    \item \textbf{Bleeding:} Significant intraoperative or postoperative hemorrhage, potentially requiring blood transfusions, reoperation for hematoma evacuation, or administration of procoagulants.
    \item \textbf{Infection:} Surgical site infection (sternal wound infection, mediastinitis), pneumonia, urinary tract infection, or systemic sepsis. Prosthetic valve endocarditis is a rare but devastating complication.
    \item \textbf{Anesthesia Complications:} Adverse reactions to anesthetic medications, respiratory depression, cardiac events during anesthesia, or prolonged intubation and ventilator dependence.
    \item \textbf{Pain:} Postoperative pain, which is typically well-managed with multimodal analgesia but can occasionally be severe or chronic.
  \end{itemize}
  \item \textbf{Procedure-Specific Risks:}
  \begin{itemize}
    \item \textbf{Stroke:} Due to embolization of plaque, air, or thrombus during surgery, potentially leading to transient ischemic attack (TIA) or permanent neurological deficits.
    \item \textbf{Heart Block:} Damage to the heart's electrical conduction system (e.g., AV node, Bundle of His) during valve excision or suture placement, potentially requiring implantation of a permanent pacemaker.
    \item \textbf{Myocardial Infarction:} Damage to the heart muscle during or after surgery, often due to coronary artery injury, inadequate myocardial protection, or coronary vasospasm.
    \item \textbf{Acute Kidney Injury (AKI):} Ranging from transient renal dysfunction to severe AKI requiring temporary or permanent dialysis, often related to cardiopulmonary bypass, hypotension, or nephrotoxic medications.
    \item \textbf{Respiratory Failure:} Prolonged mechanical ventilation, acute respiratory distress syndrome (ARDS), pleural effusions, or the need for tracheostomy.
    \item \textbf{Paravalvular Leak (PVL):} Incomplete sealing of the prosthetic valve to the native annulus, leading to regurgitation and potentially requiring reintervention (surgical or transcatheter).
    \item \textbf{Prosthetic Valve Thrombosis:} Formation of a blood clot on the prosthetic valve, leading to valve dysfunction (stenosis or regurgitation); more common with mechanical valves and inadequate anticoagulation.
    \item \textbf{Prosthetic Valve Degeneration:} Structural deterioration of bioprosthetic valves over time due to calcification or tissue fatigue, eventually requiring reoperation, particularly in younger patients.
    \item \textbf{Hemolysis:} Destruction of red blood cells, especially with mechanical valves or significant paravalvular leaks, potentially leading to anemia.
    \item \textbf{Arrhythmias:} New-onset atrial fibrillation, ventricular arrhythmias, or other dysrhythmias requiring medical management or cardioversion.
    \item \textbf{Death:} Despite significant advances in surgical techniques and perioperative care, heart valve replacement carries a small but significant mortality risk, typically ranging from 1-5\% for elective procedures, influenced by patient comorbidities and the specific valve involved.
  \end{itemize}
\end{itemize}

\section*{Postoperative Recovery Timeline}
Immediately following heart valve replacement surgery, the patient is transferred to the cardiac intensive care unit (CICU) for intensive monitoring. During the initial 24-48 hours, vital signs, cardiac rhythm, urine output, chest tube drainage, and neurological status are continuously and meticulously assessed. Pain management is promptly initiated with intravenous analgesics, and patients are typically extubated (removed from the mechanical ventilator) within hours of surgery if stable and hemodynamically optimized. Early mobilization, often within 24 hours, is strongly encouraged to prevent common postoperative complications such as deep vein thrombosis, pulmonary embolism, and pneumonia. Intravenous fluids and medications, including inotropes and vasopressors, are carefully titrated to optimize cardiac output and maintain hemodynamic stability.

Over the subsequent 1-2 weeks, patients are gradually weaned off intravenous medications, chest tubes are typically removed, and they are transferred from the CICU to a step-down unit or general ward. Physical therapy and occupational therapy commence to help patients regain strength, mobility, and independence in activities of daily living. Comprehensive education is provided regarding wound care, medication management (especially the critical importance of lifelong anticoagulation for mechanical valve recipients), and activity restrictions, particularly sternal precautions (e.g., no lifting heavy objects, no pushing/pulling for 6-8 weeks). Most patients are discharged home within 5-10 days, depending on their individual recovery trajectory and the absence of significant complications. Long-term recovery involves continued cardiac rehabilitation, strict adherence to medical regimens, and a gradual return to normal activities over several weeks to months, with full physical recovery often taking 3-6 months.

\section*{Surgical Findings \& Pathology}
During heart valve replacement, the surgeon meticulously documents the gross appearance of the excised native valve and surrounding structures. These surgical findings provide immediate insights into the etiology and severity of the valvular disease. Common observations include extensive calcification and thickening of the leaflets in aortic stenosis, myxomatous degeneration with leaflet prolapse and ruptured chordae in primary mitral regurgitation, or fusion of commissures and subvalvular apparatus in mitral stenosis. The presence of vegetations, abscesses, or annular destruction would indicate infective endocarditis. The surgeon also assesses the size of the annulus, the integrity of the surrounding myocardium, and the presence of any associated pathology.

The excised native valve tissue is routinely sent for histopathological examination. The pathology report provides a definitive microscopic diagnosis, confirming the nature of the disease. For calcific aortic stenosis, findings typically include extensive calcium deposition within the valve cusps, fibrosis, and chronic inflammatory cells. Myxomatous degeneration of the mitral valve shows accumulation of proteoglycans within the spongiosa layer, leading to leaflet thickening and elongation. Rheumatic valve disease would demonstrate chronic inflammation, fibrosis, and neovascularization. In cases of endocarditis, the pathology report would identify bacterial colonies, inflammatory infiltrates, and necrotic debris. These findings correlate with the gross surgical observations and are crucial for understanding the disease process, guiding long-term management, and informing prognosis.

\section{Expected Postoperative Course}
A normal and successful postoperative course following heart valve replacement is characterized by a progressive and steady improvement in the patient's clinical status. Initially, patients will experience incisional pain, which is effectively managed with analgesia, gradually diminishing over days to weeks. Resolution or significant improvement of preoperative symptoms such as dyspnea, angina, or syncope is expected as the new valve restores normal hemodynamics. Patients should demonstrate stable cardiac rhythm, with any new-onset arrhythmias, such as atrial fibrillation, being well-controlled medically. Wound healing should progress without signs of infection, redness, or discharge. Over time, particularly in the months following surgery, an improvement in left ventricular function and favorable cardiac remodeling (e.g., regression of hypertrophy or dilation) is anticipated, leading to enhanced exercise tolerance and overall functional capacity. For patients with mechanical valves, a stable and therapeutic international normalized ratio (INR) with appropriate anticoagulation is a critical component of a successful course.

\section{Postoperative Care \& Follow-up}
Rigorous postoperative care and long-term follow-up are essential to ensure optimal outcomes, monitor prosthetic valve function, and manage potential complications.
\begin{itemize}
  \item \textbf{Cardiology Clinic Visits:} Scheduled at 1-2 weeks post-discharge for initial wound check and medication review, then typically at 1, 3, 6, and 12 months, and annually thereafter, or more frequently if clinically indicated.
  \item \textbf{Suture/Staple Removal:} External sutures or staples are usually removed at the first post-op visit, approximately 1-2 weeks after discharge.
  \item \textbf{Echocardiography:} A baseline transthoracic echocardiogram (TTE) is performed before discharge. Repeat studies are typically done at 6-12 months and annually, or sooner if symptoms recur or there is suspicion of valve dysfunction (e.g., new murmur, signs of heart failure).
  \item \textbf{INR Monitoring:} For patients with mechanical valves, frequent INR checks are required, especially during the initial weeks, to maintain therapeutic anticoagulation levels (typically INR 2.0-3.0 or 2.5-3.5 depending on valve type and position) and prevent thrombosis or bleeding.
  \item \textbf{Blood Tests:} Routine blood work to monitor renal function, electrolytes, and complete blood counts, particularly for patients on anticoagulants, is performed periodically.
  \item \textbf{Cardiac Rehabilitation:} Often highly recommended to aid physical recovery, improve cardiovascular conditioning, and provide education on heart-healthy living, medication adherence, and risk factor modification.
  \item \textbf{Activity Resumption:} Gradual increase in physical activity is encouraged, with specific guidance on sternal precautions (e.g., lifting restrictions for 6-8 weeks) and a phased return to work and leisure activities.
  \item \textbf{Warning Signs:} Patients are educated to seek immediate medical attention for symptoms such as fever, chills, new or worsening shortness of breath, chest pain, signs of bleeding (e.g., unusual bruising, blood in urine/stool), or any neurological changes, which could indicate serious complications like infection, valve dysfunction, or stroke.
\end{itemize}
Lifelong adherence to medical therapy, including anticoagulation for mechanical valves and prophylactic antibiotics for certain procedures, is crucial.

\section*{Epidemiology}
Valvular heart disease (VHD) represents a significant global health burden, affecting millions worldwide, with its prevalence increasing substantially with advancing age. Degenerative calcific aortic stenosis (AS) is the most common indication for valve replacement in developed countries, primarily affecting older adults, with an estimated incidence of severe AS reaching 3-5\% in individuals over 75 years of age. Mitral regurgitation (MR) is also highly prevalent, with primary forms often due to myxomatous degeneration (e.g., mitral valve prolapse) and secondary (functional) forms linked to ischemic heart disease or dilated cardiomyopathy. Surgical heart valve replacement procedures are performed in hundreds of thousands of patients globally each year. While open surgical valve replacement remains a cornerstone of treatment, the advent and widespread adoption of transcatheter aortic valve replacement (TAVR) have significantly altered the epidemiological landscape, particularly for high-risk elderly patients, with TAVR procedures now outnumbering surgical AVR in many regions. Mortality rates for isolated surgical valve replacement have steadily declined over decades due to advancements in surgical techniques, myocardial protection, and perioperative care, typically ranging from 1-5\% for elective procedures, influenced by patient comorbidities and the specific valve involved. Morbidity, encompassing complications like stroke, acute kidney injury, and prolonged ventilation, also varies but has shown continuous improvement.

\section*{Outcomes \& Prognosis}
The outcomes of heart valve replacement are generally excellent, leading to significant improvements in symptoms, functional capacity, quality of life, and long-term survival for appropriately selected patients. For surgical aortic valve replacement (SAVR), 10-year survival rates typically range from 60-80\%, depending on patient age, comorbidities, and left ventricular function, with even better outcomes observed in younger, lower-risk individuals. Mitral valve replacement (MVR) also demonstrates favorable long-term survival, though operative risks can be slightly higher due to the complexity of the disease and the often more comorbid patient profile. Functional outcomes universally include resolution of heart failure symptoms, improved exercise tolerance, and often a significant reduction in left ventricular hypertrophy or dilation.

The choice between mechanical and bioprosthetic valves profoundly impacts long-term prognosis. Mechanical valves offer superior durability, often lasting the patient's lifetime, but necessitate lifelong anticoagulation with warfarin, which carries an inherent risk of bleeding (e.g., intracranial hemorrhage) and thromboembolism if INR is not therapeutic. Bioprosthetic valves, conversely, avoid the need for chronic anticoagulation (beyond the initial post-operative period) but are subject to structural valve degeneration (SVD) over time due to calcification and tissue fatigue, typically requiring reoperation within 10-15 years for younger patients, though durability is considerably longer in older individuals. Prognostic factors include patient age, preoperative left ventricular function, presence and severity of comorbidities (e.g., renal disease, pulmonary hypertension, diabetes), and the urgency of the procedure (elective vs. emergent). Recurrence of valvular dysfunction is rare with mechanical valves but is a known long-term issue for bioprosthetic valves due to SVD. Landmark clinical trials such as the PARTNER trials for TAVR have demonstrated non-inferiority or superiority of transcatheter approaches in specific patient populations, further refining prognostic considerations.

\section*{References}
\begin{enumerate}
  \item Otto CM, Bonow RO. Valvular Heart Disease. In: Zipes DP, Libby P, Bonow RO, Mann DL, Tomaselli GF, eds. Braunwald's Heart Disease: A Textbook of Cardiovascular Medicine. 12th ed. Elsevier; 2022: 1303-1360.
  \item Nishimura RA, Otto CM, Bonow RO, et al. 2020 ACC/AHA Guideline for the Management of Patients With Valvular Heart Disease: A Report of the American College of Cardiology/American Heart Association Joint Committee on Clinical Practice Guidelines. Circulation. 2020;142(14):e1-e147.
  \item MedlinePlus. Heart Valve Replacement. U.S. National Library of Medicine; 2024. Available at: \url{https://medlineplus.gov/heartvalvereplacement.html}
  \item Society of Thoracic Surgeons. Patient Information: Aortic Valve Replacement. Available at: \url{https://www.sts.org/patients/aortic-valve-replacement}
\end{enumerate}

\clearpage